% !TeX root = ../sections/02_exercises.tex

\begin{exercise}
	R-6. $A=\mathbb{Z}\oplus\mathbb{Z}$ を階数 $2$ の $\mathbb{Z}$-自由加群とし，
	その部分加群 $B,C$ を次のように置く：
	\[
	B:=\mathbb{Z}(2,0)+\mathbb{Z}(0,-3),
	\qquad
	C:=\mathbb{Z}(5,2)+\mathbb{Z}(3,5).
	\]
	\begin{enumerate}
		\item $B$ と $C$ の階数を求めよ。
		\item $A/B$ と $A/C$ の不変系を求めよ。
	\end{enumerate}
\end{exercise}

\begin{background}
	$A=\mathbb{Z}^2$ の部分加群が有限個のベクトルで生成されているとき，
	それらの生成ベクトルを列に並べた整数行列を考える。
	
	例えば，$B$ に対しては
	\[
	M_B=
	\begin{pmatrix}
		2 & 0\\
		0 & -3
	\end{pmatrix}
	\]
	を考える。
	
	部分加群の階数は，この行列の $\mathbb{Q}$ 上での階数に等しい。
	また，商加群 $A/B$ の不変系は，対応する整数行列の
	Smith 標準形を求めることで得られる。
	
	整数行列 $M$ を基本変形によって
	\[
	\operatorname{diag}(d_1,\dots,d_r)
	\]
	の形にし，
	\[
	d_1\mid d_2\mid\cdots\mid d_r
	\]
	となるようにしたものを Smith 標準形という。
\end{background}

\begin{strategy}
	$B$ については，生成ベクトル
	\[
	(2,0),\quad (0,-3)
	\]
	を列に並べた行列を考える。
	行列式は
	\[
	\det
	\begin{pmatrix}
		2 & 0\\
		0 & -3
	\end{pmatrix}
	=-6
	\]
	であるから，$B$ の階数は $2$ である。
	
	次に Smith 標準形を求める。
	\[
	\begin{pmatrix}
		2 & 0\\
		0 & -3
	\end{pmatrix}
	\]
	は Smith 標準形として
	\[
	\operatorname{diag}(1,6)
	\]
	に変形できる。
	したがって
	\[
	A/B\cong \mathbb{Z}/6\mathbb{Z}
	\]
	となる。
	
	$C$ については，生成ベクトル
	\[
	(5,2),\quad (3,5)
	\]
	を列に並べた行列
	\[
	M_C=
	\begin{pmatrix}
		5 & 3\\
		2 & 5
	\end{pmatrix}
	\]
	を考える。
	この行列式は
	\[
	5\cdot 5-3\cdot 2=19
	\]
	であるから，$C$ の階数は $2$ である。
	
	また，成分全体の最大公約数は $1$ で，行列式の絶対値は $19$ である。
	したがって Smith 標準形は
	\[
	\operatorname{diag}(1,19)
	\]
	であり，
	\[
	A/C\cong \mathbb{Z}/19\mathbb{Z}
	\]
	となる。
\end{strategy}